\documentclass[11pt,a4paper]{article}
\usepackage[utf8]{inputenc}
\usepackage[T1]{fontenc}
\usepackage{lmodern}
\usepackage[french]{babel}
\usepackage{amsmath,amssymb,mathtools}
\usepackage{icomma}
\usepackage{xcolor}
\usepackage[most]{tcolorbox}
\usepackage{enumitem}
\usepackage[a4paper,margin=2.2cm,headheight=26pt,headsep=10pt]{geometry}
\usepackage{fancyhdr}
\usepackage[hidelinks]{hyperref}
\definecolor{primary}{RGB}{222,49,99}
\definecolor{primarydark}{RGB}{150,22,55}
\definecolor{excol}{RGB}{0,135,90}
\tcbset{boxbase/.style={enhanced,breakable,boxrule=0.9pt,arc=2.5pt,
  left=3mm,right=3mm,top=2.3mm,bottom=2.3mm,fonttitle=\bfseries,coltitle=white}}
\newtcolorbox[auto counter]{corrige}[1][]{boxbase,colback=excol!4,colframe=excol,
  title={\textsc{Corrigé}~\thetcbcounter\ifx\relax#1\relax\relax\else\textmd{ --- \itshape #1}\fi}}
\newcommand{\R}{\mathbb{R}}
\newcommand{\N}{\mathbb{N}}
\newcommand{\Z}{\mathbb{Z}}
\newcommand{\ds}{\displaystyle}
\pagestyle{fancy}\fancyhf{}
\fancyhead[L]{\small\textcolor{primarydark}{\textbf{Thème 1} — Puissances, racines et exposants}}
\fancyhead[R]{\small\textcolor{primarydark}{\textsc{Corrigé — Facile}}}
\fancyfoot[C]{\small\thepage}
\renewcommand{\headrulewidth}{0.8pt}
\begin{document}

\begin{center}
{\LARGE\bfseries\color{primarydark} Niveau Facile — Corrigé}\\[2pt]
{\color{primary}\rule{0.45\textwidth}{1.2pt}}
\end{center}
\vspace{1mm}

\begin{corrige}[règles produit et quotient]
On additionne les exposants pour un produit, on les soustrait pour un quotient.
\[ \text{a) } a^{4+3}=a^{7};\quad \text{b) } a^{7-2}=a^{5};\quad \text{c) } a^{6-2}=a^{4};\quad
   \text{d) } a^{3-8}=a^{-5};\quad \text{e) } a^{2+5+1}=a^{8}. \]
\end{corrige}

\begin{corrige}[exposant nul et négatif]
Un exposant négatif signifie « inverse » ; toute base non nulle à la puissance $0$ vaut $1$.
\[ \text{a) } 5^{0}=1;\quad \text{b) } a^{-3}=\frac{1}{a^{3}};\quad \text{c) } \frac{1}{a^{4}}=a^{-4};\quad
   \text{d) } 2^{-3}=\frac{1}{2^{3}}=\frac18;\quad \text{e) } \Bigl(\tfrac13\Bigr)^{-2}=3^{2}=9. \]
\end{corrige}

\begin{corrige}[racines et exposants fractionnaires]
$\sqrt[n]{c^{p}}=c^{p/n}$.
\[ \text{a) } c^{1/5};\quad \text{b) } c^{2/3};\quad \text{c) } c^{1/4}=\sqrt[4]{c};\quad
   \text{d) } c^{6/3}=c^{2};\quad \text{e) } c^{10/2}=c^{5}. \]
\end{corrige}

\begin{corrige}[signe et parenthésage]
Le signe dans la base se neutralise pour un exposant pair ; devant, il ne participe pas à la puissance.
\begin{align*}
&\text{a) } (-3)^{2}=+9; &&\text{b) } -3^{2}=-(9)=-9; &&\text{c) } (-2)^{3}=-8;\\
&\text{d) } -2^{4}=-(16)=-16; &&\text{e) } \sqrt{(-5)^{2}}=\sqrt{25}=|-5|=5.
\end{align*}
Remarque : en a) et b) on voit bien que $(-3)^2\neq -3^2$.
\end{corrige}

\begin{corrige}[équation exponentielle à même base]
Même base $\Rightarrow$ on égale les exposants.
\begin{align*}
&\text{a) } x=5; &&\text{b) } 2x=8\Rightarrow x=4; &&\text{c) } x+1=4\Rightarrow x=3;\\
&\text{d) } 3x=0\Rightarrow x=0; &&\text{e) } x-2=3\Rightarrow x=5.
\end{align*}
Dans chaque cas $S$ est le singleton correspondant, par ex. a) $S=\{5\}$.
\end{corrige}

\end{document}
